\documentclass[11pt,a4paper]{article}
\usepackage{gtu-el-yazisi}
\begin{document}
\notbilgi{RF Tasarım}{Ders 1: RF'ye giriş ve temel kavramlar}{Ders1.pdf}

\sayfa{1}
Selami Şahin\par
\textit{(Hocanın telefon ve e-posta bilgileri aktarılmadı.)}

\begin{itemize}[label={},leftmargin=1.5em,itemsep=2pt]
\item 1- RF Basic Concepts \& components
\item $\left.\begin{tabular}[c]{@{}l@{}}2- RF design Theory \& principle\\ 3- RF design Architecture (Rx \& Tx)\end{tabular}\right\}$ \textit{\okunamadi{4} özellikle en iyi anlamamız gerekiyor}
\item 4- Modulation in Comm. Sys.
\item 5- RF Microwave \& Transmission line
\item 6- RF \& Microwave Filters
\item 7- RF Networks \quad (S parametreleri, saçılma \ldots)
\item 8- RF components \quad (Mixers, switches, Oscilator, \ldots)
\end{itemize}

\begin{itemize}[label={},leftmargin=1.5em,itemsep=2pt]
\item \%30 Arasınav
\item \%10 ödev (4-5 adet)
\item \%40 final
\item \%20 Project $\rightarrow$ Bir tane makale seçiyorsunuz, ve bu makaleyi sunuyorsunuz, sunum yapıyorsunuz.
\end{itemize}

\elcizimi[0.35]{sekiller/Ders1-s1-1.png}[1]

\[
\underbrace{\overbrace{C}^{\text{kanal kapasitesi}}}_{\text{bit/sn}} = \overbrace{W}^{\text{Hz (kanalın genişliği)}} \log_2\Bigl(1+\underbrace{\mathrm{SNR}}_{\text{signal to noise ratio}}\Bigr)
\]
\textit{Shannon's Capacity Theorem}

\elcizimi[0.75]{sekiller/Ders1-s1-2.png}[1]

\elcizimi[0.8]{sekiller/Ders1-s1-3.png}[1]

Televizyon \qty{10}{\giga\hertz}

ITU kurumu nerde haberleşeceğini söyler

\sayfa{2}
RF ısıma yapar, anten çalıştığı bölgeye göre tasarlanmış olacak.

\elcizimi[0.7]{sekiller/Ders1-s2-1.png}[2]

Gürültüyü minimalize etmeliyiz

RF: 3 kHz frekansından 300 GHz'e kadar olan tüm sinyaller. (100 GHz üstü milimitre wave, 10 GHz üstü microwave)

\elcizimi[0.45]{sekiller/Ders1-s2-2.png}[2]

{\renewcommand{\arraystretch}{1.3}
\begin{tabular}{@{}l@{\hspace{1.5em}}l@{\hspace{1.5em}}l@{\hspace{1em}}l@{}}
 & $f$ & $\lambda = c/f$ & \\
\midrule
\begin{tabular}[t]{@{}l@{}}VLF:\\ {\small (Very low frequency)}\end{tabular} & 3 -- 30 kHz & 33\footnotemark\ km $\sim$ 10 km & \\
LF: & 30 $\sim$ 300k & 10 km $\sim$ 1 km & \\
MF: & 300k $\sim$ 3M & 1 km $\sim$ 100 m & \\
HF: & 3M $\sim$ 30M & 100 m $\sim$ 10 m & \\
VHF: & 30M $\sim$ 300M & 10 m $\sim$ 1 m & $\rightarrow$ TV, FM \\
UHF: & 300M $\sim$ 3 GHz & 1 m $\sim$ 100 mm & $\rightarrow$ TV, mobile, GPS, WiFi, 4G \\
SHF: & 3G $\sim$ 30 GHz & 100 mm $\sim$ 10 mm & $\rightarrow$ WiFi, satellite \\
EHF: & 30G $\sim$ 300 GHz & 10 mm $\sim$ 1 mm & $\rightarrow$ Astronomy, satellite \\
\end{tabular}}
\footnotetext{Claude notu: muhtemelen 100 km olmalı.}

\elcizimi[0.5]{sekiller/Ders1-s2-3.png}[2]

\elcizimi[0.35]{sekiller/Ders1-s2-4.png}[2]
frekans büyüdükçe anten küçülür.

\section{Noise}
\elcizimi[0.7]{sekiller/Ders1-s2-5.png}[2]

\sayfa{3}
Varyans güç

\elcizimi[0.5]{sekiller/Ders1-s3-1.png}[3]
\elcizimi[0.45]{sekiller/Ders1-s3-2.png}[3]
\elcizimi[0.5]{sekiller/Ders1-s3-3.png}[3]

\subsection{White Noise}
\elcizimi[0.5]{sekiller/Ders1-s3-4.png}[3]

\subsection{Pink noise}
\elcizimi[0.5]{sekiller/Ders1-s3-5.png}[3]

White Noise $\sim$ Thermal Noise (Shot Noise) \quad [$\rightarrow$ White Noise $N(m_0,\sigma)$]

Flicker Noise ($\frac{1}{f^{n}}$ noise) $\sim$ Pink noise \\
(Low frequency)

+ üst üste çizersek

\elcizimi[0.55]{sekiller/Ders1-s3-6.png}[3]

SNR oranı

\elcizimi[0.7]{sekiller/Ders1-s3-7.png}[3]

\sayfa{4}
\[
\mathrm{SNR} = \frac{P_s[\mathrm{W}]}{P_n[\mathrm{W}]} \qquad \overset{\text{birimsiz}}{[\,\cdot\,]}
\]

\begin{align*}
\mathrm{SNR}_{[\mathrm{dB}]} &= 10\log\left[\frac{P_s}{P_n}\right]\mathrm{dB} && \text{\shortstack[l]{sadece logaritma birimi bel\\ 10 ile çarparsanız decibel}}\\
& \phantom{=}\ \log[\mathrm{SNR}] \\
10\log[2\cdot\mathrm{SNR}] &= \underbrace{10\log 2}_{3\,\mathrm{dB}} + 10\log[\mathrm{SNR}]
\end{align*}

2 kat arttırırsak 3 db \\
10 kat 10 db \quad \okunamadi{} \\
100 kat 20 db

\[
\begin{array}{@{}l@{\quad}l@{\quad}l@{}}
10\ \mathrm{W} & \rightarrow & \mathrm{dBW}\\
10\log\left[\dfrac{P_s}{1\,\mathrm{W}}\right] & \Rightarrow & \mathrm{dBW}\\
10\log\left[\dfrac{P_s}{1\,\mathrm{mW}}\right] & \Rightarrow & \mathrm{dBm}
\end{array}
\]

\[
10\log[P_s] + 10\log\left[\frac{1}{1\,\mathrm{mW}}\right] = \bigl(P_s[\mathrm{dBW}] + 30\ \mathrm{dB}\bigr)\,[\mathrm{dBm}]
\]

\begin{align*}
10\log\left[\frac{1\,\mathrm{mW}}{1\,\mathrm{mW}}\right] &= 0\ \mathrm{dBm}\\
10\log\left[\frac{10\,\mathrm{mW}}{1\,\mathrm{mW}}\right] &= 10\ \mathrm{dBm}
\end{align*}

\elcizimi[0.5]{sekiller/Ders1-s4-1.png}[4]

\elcizimi[0.85]{sekiller/Ders1-s4-2.png}[4]

\sayfa{5}
\elcizimi[0.9]{sekiller/Ders1-s5-1.png}[5]

\elcizimi[0.35]{sekiller/Ders1-s5-2.png}[5]

RF Front-End $=$ RFE

\elcizimi[0.7]{sekiller/Ders1-s5-3.png}[5]

Aynı anda hem verici hem alıcı çalışmaz

Half duplex

\elcizimi[0.5]{sekiller/Ders1-s5-4.png}[5]

Full duplex

\elcizimi[0.35]{sekiller/Ders1-s5-5.png}[5]

Simplex

\sayfa{6}
\section{Antennas}
Log-periodic Antenna \quad her iki \okunamadi{yönden} alır

\elcizimi[0.5]{sekiller/Ders1-s6-1.png}[6]

metal reflektör koyarsanız sadece tek yönlü $\leftarrow$ olur

\elcizimi[0.8]{sekiller/Ders1-s6-2.png}[6]

\elcizimi[0.65]{sekiller/Ders1-s6-3.png}[6]

\section{Patch Antenna}
\elcizimi[0.55]{sekiller/Ders1-s6-4.png}[6]

\elcizimi[0.65]{sekiller/Ders1-s6-5.png}[6]

\sayfa{7}
\section{RF Filters}
\[ 20\log\frac{V_o}{V_i} = 10\log\left[\frac{P_o}{P_i}\right] = G \]

\elcizimi[0.7]{sekiller/Ders1-s7-1.png}[7]

\elcizimi[0.55]{sekiller/Ders1-s7-2.png}[7]
\[ BW = f_2 - f_1 \]
\elcizimi[0.6]{sekiller/Ders1-s7-3.png}[7]

\elcizimi[0.45]{sekiller/Ders1-s7-4.png}[7]

\elcizimi[0.55]{sekiller/Ders1-s7-5.png}[7]
BSP

Aktif komponentler akımı yön veriyor. \okunamadi{}

\section{LNA}
\elcizimi[0.35]{sekiller/Ders1-s7-6.png}[7]

\elcizimi[0.4]{sekiller/Ders1-s7-7.png}[7]
\[ \mathrm{SNR}_{in} = \frac{S_i}{N_i} \qquad \mathrm{SNR}_{out} = \frac{S_o}{N_o+N_A} = \frac{S_i\cdot G}{N_i\,G+N_A} \]
$N_A <<<$ seçilmeli \quad $\swarrow$ Filtrenin kendisinden kaynaklanan ek gürültü

\elcizimi[0.4]{sekiller/Ders1-s7-8.png}[7]

\[ \mathrm{SNR}_{out} < \mathrm{SNR}_{in} \]

\section{Requirements}
\begin{itemize}[label=--,leftmargin=1.5em]
\item Low noise coming from device ($\Rightarrow$ low NF) \\ $\hookrightarrow$ Noise Figure
\item Linearity \quad ($\Rightarrow$ high $P_{1\mathrm{dB}}$ $\rightarrow$ 1dB Compression Point \\ \hspace*{5em} IP3 $\rightarrow$ Third order Intercept Point)
\item Low Power Consumption
\item Low Area \quad $''$ \quad (Small form factor)
\item Gain ($\Rightarrow S_{21}$)
\item Matching to anten filter
\end{itemize}

\end{document}
